\chapter*{Заключение}
\addcontentsline{toc}{chapter}{Заключение}

В рамках настоящего исследования была рассмотрена задача разработки сервиса коллективных музыкальных рекомендаций для заведений и мероприятий. Работа объединила анализ предметной области и рынка музыкального сопровождения, правовую и продуктовую проработку возможного сервиса, а также техническое исследование методов формирования групповых музыкальных рекомендаций на основе открытого датасета YAMBDA.

Проведенный обзор литературы, посвященной музыкальному сопровождению бизнеса и сервисных пространств, показал, что музыка является значимым элементом клиентского опыта и может влиять на эмоциональные реакции, восприятие атмосферы, удовлетворенность и поведенческие намерения посетителей. При этом наиболее устойчивые положительные эффекты связаны не с самим фактом наличия музыки, а с ее субъективной привлекательностью для слушателя и соответствием контексту пространства. Следовательно, адаптация музыкального потока под фактический состав аудитории может рассматриваться как перспективный способ повысить вероятность того, что звучащая музыка окажется релевантной и приятной для конкретных посетителей. В настоящей работе такая адаптация предложена через механизм групповых рекомендаций, использующих накопленные музыкальные предпочтения пользователей стриминговой платформы.

Анализ существующих практик музыкального сопровождения показал, что значительная часть сервисов музыки для бизнеса фактически закрывает для клиента задачу недорогой легализации фоновой музыки и снижения риска претензий со стороны организаций по коллективному управлению правами. Однако такие решения часто опираются на ограниченные каталоги или заранее подготовленные плейлисты и потому не всегда позволяют использовать популярную музыку, ожидаемую посетителями в барах, клубах, на танцевальных мероприятиях и иных музыкально чувствительных площадках. Другие сервисы, включая решения по заказу музыки посетителями, допускают работу с популярными треками и могут предусматривать отчетность для РАО и ВОИС, но при этом проблема легального источника фонограмм не всегда раскрывается достаточно ясно. На практике воспроизведение популярной музыки нередко сводится к использованию личных аккаунтов стриминговых сервисов, локальных аудиофайлов или иных смешанных схем. В событийном сегменте также встречаются прямые разрешения правообладателей на публичное исполнение произведений в рамках конкретного мероприятия, включая практику так называемых <<отказных>> для РАО и ВОИС.

С учетом этой правовой неопределенности в работе сделан вывод, что проектируемый сервис не должен рассматриваться как простая замена лицензионного B2B-плеера. Более корректной является модель информационно-рекомендательной надстройки, которая формирует музыкальные рекомендации и может работать поверх уже урегулированного контура публичного воспроизведения. В качестве возможных режимов рассмотрены использование рекомендаций в сервисных нуждах бизнес-аккаунта, сопоставление собственных аудиофайлов клиента с библиотекой платформы через цифровые отпечатки, метаданные и аудиоэмбеддинги, а также интеграция с существующими легальными B2B-решениями. Для минимизации юридических и репутационных рисков доступ к функциональности сервиса должен быть связан с наличием у бизнеса договоров с РАО и ВОИС либо иного правового основания для публичного воспроизведения, а сам сервис должен автоматически формировать отчетность о фактически использованных произведениях и фонограммах.

Во второй главе на этой основе была разработана концепция сервиса, включающая B2B-направление <<Яндекс Музыка для бизнеса>> и B2C-функцию <<Резонанс>> для пользователей Яндекс Музыки. Пользовательское приложение в предложенной модели фиксирует идентификаторы BLE-маяков заведения или мероприятия и уведомляет пользователя о возможности добровольно повлиять на музыкальную среду через сценарий <<Резонировать>> или <<Войти в резонанс>>. Через этот интерфейс пользователь получает объяснение механики и дает согласие на обработку данных для формирования коллективной рекомендации. Серверный контур получает идентификатор пользователя и идентификатор маяка, сопоставляет их с конкретным бизнес-объектом, формирует текущую эфемерную группу посетителей и рассчитывает рекомендательный поток для бизнес-аккаунта организации. При этом заведению передается только результат в виде списка рекомендуемых треков или очереди воспроизведения; индивидуальные музыкальные профили, история прослушиваний и персональные данные посетителей остаются внутри платформенного контура.

Важным условием применимости сервиса является сохранение контроля со стороны бизнеса. Поэтому концепция предполагает наличие фильтров и настроек по жанрам, уровню энергичности, времени суток, стоп-листам, допустимой популярности, языку и другим параметрам музыкальной среды. Такие настройки могут задаваться классическим интерфейсом, текстовым описанием с участием AI-консультанта или через использование истории прослушиваний выбранного аккаунта, которому может быть придан повышенный вес как условному музыкальному профилю заведения. Тем самым сервис не подменяет музыкальную политику площадки, а работает внутри заданных ею рамок.

В качестве приоритетных сегментов для первых пилотных запусков были выделены бары, рестораны и клубы, проводящие танцевальные мероприятия, concert- и event-площадки, а также другие форматы, предполагающие использование популярной музыки и наличие лицензионного контура взаимодействия с РАО и ВОИС. На основе сценарной оценки рынка было показано, что наиболее релевантным является не весь массив заведений, где звучит музыка, а сегмент музыкально чувствительных точек. В текущей версии работы TAM такого рынка оценен в 14--23 тыс. точек, SAM~--- в 3--5 тыс. точек, а реалистично достижимый SOM для раннего коммерческого масштабирования~--- в 450--550 бизнесов без механического включения событийного формата. Рекомендуемая ценовая гипотеза подписки в 6\,000--9\,000 рублей в месяц за заведение позволяет рассматривать MVP и пилотный запуск как экономически проверяемые при условии использования существующей инфраструктуры экосистемы Яндекса.

Целевыми аудиториями сервиса определены владельцы и арт-директора заведений, организаторы мероприятий, DJ, пользователи Яндекс Музыки, а также сама компания Яндекс как экосистемный бенефициар. Для бизнеса ценность заключается в управляемой адаптации музыкальной атмосферы, снижении административной нагрузки и возможности автоматизированной отчетности. Для организаторов и DJ сервис может стать дополнительным сигналом о предпочтениях аудитории, особенно в сценариях разогрева и поддержки вовлеченности. Для пользователя ценность состоит в мягком влиянии на музыку без прямого заказа треков и без раскрытия персонального профиля заведению. Для Яндекса подобный сервис может создать уникальное конкурентное отличие Яндекс Музыки и сформировать новый класс данных о групповом и пространственном музыкальном потреблении. При этом такие данные могут использоваться только в агрегированном, обезличенном и правомерном виде.

В третьей главе были систематизированы теоретические и технологические основы построения музыкальных рекомендательных систем. Рассмотрение классической коллаборативной фильтрации, последовательных моделей и трансформерных архитектур SASRec, BERT4Rec и gSASRec позволило обосновать выбор SASRec в качестве персонального последовательного скорера. Особое значение для работы имеет специфика музыкального домена: пользователи регулярно повторно слушают уже знакомые треки, поэтому стандартные протоколы оценки, заимствованные из рекомендаций фильмов или товаров, не всегда применимы к музыке. Также было показано, что аудиоэмбеддинги являются важным содержательным сигналом, поскольку они доступны независимо от частоты прослушивания трека и могут использоваться для построения пользовательских аудиопрофилей.

В четвертой главе были рассмотрены подходы к формированию коллективных рекомендаций. Были описаны классические стратегии агрегации индивидуальных оценок, включая Average, Least Misery и Max Pleasure, а также различие между устойчивыми и эфемерными группами. Для настоящей работы ключевым является именно эфемерный сценарий, поскольку посетители заведения или мероприятия образуют временную группу без собственной истории совместного потребления. В качестве релевантных нейросетевых бейзлайнов были выбраны AGREE и GroupIM, а AlignGroup был рассмотрен как методологически важная, но неприменимая в данной постановке архитектура из-за зависимости от persistent-групп и совместного обучения с индивидуальной веткой. На этой основе была сформулирована центральная техническая гипотеза: замена обучаемых ID-эмбеддингов на аудиосигнал в attention-агрегаторе должна повысить качество групповой музыкальной рекомендации.

В пятой главе была реализована и экспериментально исследована модель коллективных музыкальных рекомендаций на датасете YAMBDA-50m. На персональном уровне был воспроизведен яндексовский SASRec-бейзлайн: итоговое значение test NDCG@10 составило 0{,}0726 против 0{,}0742 в официальном ориентире. В ходе воспроизведения был получен самостоятельный методологический результат: в музыкальном домене нельзя механически маскировать историю пользователя при ранжировании, поскольку значительная часть релевантных треков относится к повторному потреблению. Отказ от такого маскирования дал основной прирост качества и позволил закрыть setup-gap между исходной реализацией и baseline.

На групповом уровне была использована двухэтапная схема: замороженный персональный скорер формирует индивидуальные оценки треков, после чего обучаемый агрегатор преобразует их в групповую оценку. В эксперименте сравнивались три классические стратегии агрегации, два ID-based attention-агрегатора и два предложенных audio-based attention-агрегатора: Audio-AGREE и Group Cross-Attention. Результаты показали, что audio-aware методы устойчиво превосходят ID-based методы по NDCG@10 и NDCG@20. Paired bootstrap подтвердил положительную разность во всех восьми сравнениях audio- и ID-методов, а абсолютный выигрыш составил примерно 0{,}009--0{,}013 NDCG, что соответствует относительному приросту около 12--16\% от уровня ID-агрегаторов. Дополнительно было показано, что преимущество audio-методов особенно заметно на группах большего размера, где возрастает сложность согласования разных музыкальных предпочтений.

Полученный результат подтверждает основную техническую гипотезу исследования. ID-based attention поверх замороженного персонального скорера оказался ограничен структурным потолком: AGREE и GroupIM практически совпадают по качеству и не дают существенного преимущества над сильной тривиальной стратегией Max Pleasure. Напротив, аудиоэмбеддинги предоставляют независимый содержательный сигнал о музыкальной близости трека и пользовательского аудиопрофиля, что позволяет улучшить групповое ранжирование. Таким образом, использование audio-aware attention является обоснованным направлением для дальнейшего развития групповых музыкальных рекомендаций.

В шестой главе экспериментальная модель была связана с демонстрационным сценарием применения. Алгоритмический прототип показывает, как список идентификаторов пользователей из YAMBDA задает текущую эфемерную группу, для которой формируется кандидатное множество на основе top-200 персональных кандидатов каждого участника и строится ранжированный список анонимизированных треков. Веб-демонстрация концептуализирует динамическое изменение состава группы и пересчет рекомендаций при появлении или исчезновении пользователей. Отдельно был рассмотрен Android BLE-прототип, который позволяет сохранять BLE-маяки по MAC-адресам, сканировать окружающие Bluetooth-устройства и выбирать ближайший сохраненный маяк по наиболее сильному свежему RSSI-сигналу. Данный компонент не является системой точной геолокации, но иллюстрирует техническую возможность использовать BLE как приближенный индикатор принадлежности пользователя к зоне пространства.

Работа имеет ряд ограничений. Во-первых, исследование не использует внутренние данные, промышленные алгоритмы и B2B-лицензионную инфраструктуру Яндекс Музыки, поэтому сервисная часть описана как концепция, а не как готовая промышленная реализация. Во-вторых, датасет YAMBDA обезличен: в нем отсутствуют названия треков, исполнители, жанры и человекочитаемые пользовательские профили, поэтому качество модели оценивается метриками ранжирования, а не субъективным восприятием реального плейлиста. В-третьих, группы в эксперименте синтезируются из индивидуальных профилей, а не наблюдаются в реальных заведениях или на мероприятиях. В-четвертых, промышленная реализация потребует отдельной проверки пользовательских согласий, правовой модели публичного воспроизведения, BLE-инфраструктуры, интерфейсов бизнес-профиля и устойчивости сервиса в реальных условиях.

Перспективы дальнейшего развития включают несколько направлений. На продуктовой стороне необходимы пилотные запуски с заведениями, организаторами мероприятий, DJ и пользователями, позволяющие проверить готовность к использованию сервиса и уточнить ценностное предложение. На правовой стороне требуется детальная проработка B2B-лицензирования, отчетности и модели обработки персональных данных. На технической стороне перспективны фильтрация и бизнес-ограничения при формировании плейлиста, неравномерное взвешивание вклада отдельных пользователей, учет нескольких BLE-зон, fairness-метрики для групп, cold-start-срезы, end-to-end дообучение персонального скорера под групповую задачу и масштабирование экспериментов на более крупные версии YAMBDA. Отдельными направлениями могут стать DJ-ассистент и B2B-аналитика на основе агрегированных данных о составе аудитории, музыкальном потоке и пространственном поведении посетителей.

Таким образом, в работе показано, что сервис коллективных музыкальных рекомендаций для заведений и мероприятий может занять содержательную продуктовую нишу между классическими B2B-сервисами фоновой музыки и сервисами индивидуального заказа треков. Его ценность определяется сочетанием трех условий: учет фактического состава аудитории, сохранение управляемости музыкальной атмосферы для бизнеса и соблюдение правовых и приватностных ограничений. Экспериментальная часть на YAMBDA демонстрирует, что техническая основа такого сервиса реализуема: audio-aware агрегаторы позволяют улучшить качество группового ранжирования относительно ID-based подходов, а значит, дальнейшая проверка концепции в реальных пилотных условиях представляется обоснованной.
